\chapter{Modules, Packages, and Robust Exception Handling Architecture}
\label{chap:modules-exceptions}

\section{Modular Software Architecture in Python}
A \textbf{module} is a single Python source file (\texttt{.py}). A \textbf{package} is a filesystem directory containing an \texttt{\_\_init\_\_.py} file that establishes a hierarchical namespace.

\section{Module Loading and Execution Mechanics}
\begin{lstlisting}
# Standard entry point execution guard
if __name__ == "__main__":
    # Code executed only when script is launched directly from command line
    # NOT executed when imported by external modules
    print("Executing as standalone script.")
\end{lstlisting}

\section{Exception Propagation and Stack Unwinding}
\begin{definition}[Exception Unwinding]
When an exception is raised, the PVM halts linear instruction flow and traverses the call stack upward. If an enclosing \texttt{try-except} block matches the exception type, execution resumes inside that block. If the stack unwinds completely without a handler, CPython prints a \texttt{Traceback} to \texttt{sys.stderr} and terminates with non-zero exit code.
\end{definition}

\section{The Complete \texttt{try-except-else-finally} Pipeline}
\begin{lstlisting}
try:
    file_handle = open("telemetry.dat", "r")
    payload = file_handle.read()
except FileNotFoundError as err:
    print(f"I/O Failure: {err}")
else:
    # Executed ONLY if NO exception was raised in try block
    print(f"Successfully ingested {len(payload)} bytes.")
finally:
    # Executed UNCONDITIONALLY (clean up resources regardless of success or crash)
    print("Closing file descriptors.")
\end{lstlisting}

\section{Custom Exception Hierarchies}
\begin{lstlisting}
class DataPipelineError(Exception):
    '''Base exception for data processing pipeline failures.'''
    pass

class SchemaValidationError(DataPipelineError):
    def __init__(self, missing_fields, message="Schema mismatch detected"):
        super().__init__(f"{message}: {missing_fields}")
        self.missing_fields = missing_fields

class NetworkTimeoutError(DataPipelineError):
    pass
\end{lstlisting}

\section{Context Managers and the \texttt{with} Statement}
\begin{keyequation}[Context Manager Protocol (PEP 343)]
Any class implementing the context protocol satisfies:
\begin{equation}
\text{with Context() as x:} \iff \begin{cases}
x = \text{Context}.\_\_\text{enter}\_\_() \\
\text{try: body} \\
\text{finally: Context}.\_\_\text{exit}\_\_(\text{exc\_type, exc\_val, exc\_tb})
\end{cases}
\end{equation}
\end{keyequation}

\section{Chapter Summary}
\begin{summarybox}
$\bullet$ The \texttt{sys.path} array dictates the resolution order of imported modules.\\
$\bullet$ Always catch specific exceptions rather than generic \texttt{Exception} or bare \texttt{except:}.\\
$\bullet$ The \texttt{finally} block is guaranteed to execute even across \texttt{return} statements.
\end{summarybox}
